% Using InitechBase, Chapter 3: About Program Files (annex Sec 8.2 U3).
% Sources: samir_main.c repl_try_do_file / repl_load_prg (DO <name> loads
% <name>.PRG, 8191-byte cap); cmd/proc.c proc_run (main program up to the
% first PROCEDURE; DO of a procedure of the same source; scope model);
% cmd/flow.c (one command per line, comments, structures, 512 lines); the
% REPL-owned verbs USE/CLOSE/QUIT/EXIT and DO <file> are dot prompt only
% (observed: in a program they are #16).
\chapter{About Program Files}
\Gates{test-pubs-examples, test-interp-proc, test-interp-flow}
A \term{program} is a list of InitechBase commands kept in a file.  When you
type \cmd{DO} and the name of the program at the dot prompt, InitechBase reads
the file and carries out its commands one after another, as if you had typed
them, until the program ends.  The work you do often -- a month-end report, a
search, a correction of many records -- is done most surely by a program.

\section{Writing a Program File}
A program file is a text file with the extension \file{.PRG}.  Write it with
any editor that saves plain text, such as the InitechDOS line editor
\file{EDLIN} (\doctitle{InitechDOS Reference}).

Each line of the file holds one command, of up to 254 characters.  A program
file may hold up to 512 lines and 8,191 characters.  Blank lines are allowed,
and so is indentation: InitechBase ignores the blanks at the beginning of a
line, so indent the commands inside \cmd{IF}, \cmd{DO WHILE} and
\cmd{DO CASE} to show the structure of the program.  A line that begins with
an asterisk or \cmd{NOTE} is a comment, and two ampersands begin a comment at
the end of a command.

\section{The Commands of a Program}
A program may use every command of Chapter~5 except those that belong to the
dot prompt: \cmd{USE}, \cmd{CLOSE}, \cmd{QUIT} and \cmd{EXIT}, and \cmd{DO} of
another program file.  A program works on the databases that are open when it
is run: open them at the dot prompt with \cmd{USE}, and select the work area
the program expects, before you type \cmd{DO}.  The program may select other
work areas with \cmd{SELECT}.

Three structures exist only in programs: \cmd{IF} \ldots{} \cmd{ENDIF}
chooses between two groups of commands, \cmd{DO CASE} \ldots{} \cmd{ENDCASE}
among several, and \cmd{DO WHILE} \ldots{} \cmd{ENDDO} repeats a group.  Each
condition they test must be a logical expression.  \cmd{ACCEPT}, \cmd{INPUT}
and \cmd{WAIT} ask the user for information while the program runs; \cmd{?}
and \cmd{??} display results.

\section{Running a Program}
Type \cmd{DO} and the name of the program file; the extension \file{.PRG} may
be left out:
\DotExample{examples/u3-run}
The program runs from its first line.  It ends at a \cmd{RETURN}, at its first
\cmd{PROCEDURE} line or after its last line, and InitechBase displays the dot
prompt.  The memory variables the main program created remain after it ends,
as if they had been created at the dot prompt; those created by its
procedures are released when the procedures return.

When a command of the program cannot be carried out, the program stops at
that command.  InitechBase displays the message of the error (Chapter~7) and
the dot prompt.  The message follows the last output of the program on the
same line when that output did not end its line.  An \cmd{ON ERROR} command
given before the program runs is carried out first.

\section{Procedures}
A \term{procedure} is a part of a program file that the program calls by name
with \cmd{DO}.  The procedures follow the main program in the same file.  Each
begins with a \cmd{PROCEDURE} line and ends with \cmd{RETURN}.  When a
procedure returns, the program goes on after the \cmd{DO} that called it.

Values are passed to a procedure with \cmd{DO} \ldots{} \kw{WITH} and received
with \cmd{PARAMETERS}.  A procedure that must give back a result stores it in a
memory variable declared \cmd{PUBLIC}, as \file{RAISE.PRG} does:
\ProgramListing{RAISE.PRG}
\DotExample{examples/u3-proc}

\section{Memory Variables in Programs}
A memory variable created by a procedure is \term{private} to it: the
procedures it calls can use and change it, but it is released when the
procedure that created it returns.  \cmd{PRIVATE} creates private
memory variables explicitly, hiding any of the same names that the callers
have; \cmd{PUBLIC} creates memory variables that every program can use and
that remain after the program ends.  \cmd{RELEASE} removes memory variables at
once.

\section{The Programs on the Sample Diskette}
The Sample Programs and Utilities diskette holds the programs used as examples
in this manual: \file{STAFFLST.PRG}, \file{PAYTOTAL.PRG}, \file{DEPTNAME.PRG},
\file{RAISE.PRG} and \file{TALLY.PRG} with \file{STAFF.DBF}, and the programs of
the Initech Accounting System with \file{INVOICE.DBF}, which
\doctitle{Programming with InitechBase} describes in full.  Each begins with
comment lines that name the database it expects.
